\documentclass[11pt]{article}
\usepackage[margin=1in]{geometry}
\usepackage{amsmath,amssymb,amsthm}
\usepackage{graphicx}
\usepackage[colorlinks=true,linkcolor=blue,citecolor=blue,urlcolor=blue]{hyperref}

\newtheorem{theorem}{Theorem}
\newtheorem{lemma}{Lemma}
\newtheorem{remark}{Remark}

\newcommand{\N}{\mathcal N}
\newcommand{\om}{\ominus}

\title{A counterexample to the N-function factorization question}
\author{Agent lane 8}
\date{July 3, 2026}

\begin{document}
\maketitle

\section*{Status}

This packet gives a candidate counterexample to Problem 6.a.1 in Kiwerski and
Tomaszewski, \emph{A few last words on pointwise multipliers of
Calderon-Lozanovskii spaces} \cite{KT}. The verdict is likely valid, subject to
human review of the piecewise-linear inverse construction.

\section*{Source problem}

The source paper asks whether the usual inverse relation for factorization is
automatic once all three Orlicz functions involved are N-functions.

\begin{center}
\includegraphics[width=0.98\linewidth]{figures/open_problem_crop.png}
\end{center}

In the notation of the source, the generalized Young conjugate is
\[
  (G\om F)(t)=\sup_{s\geq 0}\{G(st)-F(s)\}.
\]
Problem 6.a.1 asks whether
\[
  F^{-1}(G\om F)^{-1}\approx G^{-1}
\]
must hold whenever $F$, $G$, and $G\om F$ are N-functions.

\section*{Proof intuition}

Take $G(t)=t^2$. If $H=G\om F$, then $H(t)\leq u$ forces
\[
  t\leq \inf_{s>0}\frac{\sqrt{F(s)+u}}{s}.
\]
Thus $H^{-1}(u)$ can be made small by finding a large $s$ for which $F(s)$ is
much larger than $u$ but $s$ is even larger relative to $F^{-1}(u)$.

The construction builds $F$ through its inverse $A=F^{-1}$. Convexity of $F$ is
equivalent to concavity of $A$. We choose $A$ concave, increasing, unbounded,
and with $A(u)/\sqrt u\to0$, so $F$ grows faster than $t^2$ and $H=G\om F$ is
finite and again an N-function. But $A$ is arranged to have arbitrarily long
almost-linear stretches:
\[
  A(u_n y_n)\geq c\,y_n A(u_n),\qquad y_n\to\infty.
\]
Evaluating the above estimate at $s=A(u_ny_n)$ gives
\[
  \frac{F^{-1}(u_n)H^{-1}(u_n)}{G^{-1}(u_n)}
  \lesssim \frac{\sqrt{y_n}}{y_n}\to0.
\]
This destroys the lower half of the desired equivalence.

\begin{lemma}[A concave inverse with long linear stretches]\label{lem:A}
There is a continuous, strictly increasing, concave, unbounded function
$A:[0,\infty)\to[0,\infty)$ such that $A(0)=0$,
\[
  \frac{A(u)}{u}\to\infty\quad(u\downarrow0),\qquad
  \frac{A(u)}{u}\to0,\quad
  \frac{A(u)}{\sqrt u}\to0\quad(u\to\infty),
\]
and there are numbers $u_n\to\infty$ and $y_n\to\infty$ with
\[
  A(u_ny_n)\geq \frac{y_n}{4}A(u_n).
\]
\end{lemma}

\begin{proof}
Put $A(u)=u^{1/4}$ on $[0,1]$. We extend $A$ piecewise linearly on
$[1,\infty)$ with strictly decreasing positive slopes. Let $v_0=1$ and let
$\alpha_0=1/4$ be the left derivative at $1$.

Assume that $A$ has been defined up to $v_{n-1}$, with terminal slope
$\alpha_{n-1}>0$, and put $c_{n-1}=A(v_{n-1})$. Choose
$y_n=\max\{2,n^4\}$. Pick $u_n>v_{n-1}$ so large that, with
\[
  m_n=\frac{1}{n\sqrt{u_ny_n}},
\]
we have
\[
  m_n<\alpha_{n-1}/2,\qquad
  c_{n-1}\leq m_nu_n,\qquad
  m_nu_n(y_n-1)>n.
\]
This is possible because $m_n\to0$ while
$m_nu_n=\sqrt{u_n}/(n\sqrt{y_n})\to\infty$ as $u_n\to\infty$. Define
$v_n=u_ny_n$ and extend $A$ linearly with slope $m_n$ on $[v_{n-1},v_n]$.
Let $\alpha_n=m_n$.

The slopes decrease, so the resulting function is concave and strictly
increasing. It is unbounded because
$A(v_n)\geq A(u_n)+m_nu_n(y_n-1)>n$. Also,
\[
  A(v_n)=c_{n-1}+m_n(v_n-v_{n-1})
  \leq m_nu_n+m_nv_n\leq 2m_nv_n,
\]
and hence
\[
  \frac{A(v_n)}{\sqrt{v_n}}\leq 2m_n\sqrt{v_n}=\frac{2}{n}.
\]
On the whole interval $[v_{n-1},v_n]$,
\[
  \frac{A(u)}{\sqrt u}
  \leq \frac{A(v_{n-1})}{\sqrt{v_{n-1}}}+m_n\sqrt{v_n}.
\]
The first term is at most $2/(n-1)$ for $n\geq2$, while the second term is
$1/n$. Therefore $A(u)/\sqrt u\to0$. This implies $A(u)/u\to0$ at infinity.
At zero, $A(u)/u=u^{-3/4}\to\infty$.

Finally,
\[
  A(u_n)\leq c_{n-1}+m_nu_n\leq 2m_nu_n,\qquad
  A(u_ny_n)\geq m_n(u_ny_n-u_n).
\]
Since $y_n\geq2$,
\[
  \frac{A(u_ny_n)}{A(u_n)}
  \geq \frac{m_nu_n(y_n-1)}{2m_nu_n}
  \geq \frac{y_n}{4}.
\]
\end{proof}

\begin{theorem}
There exist N-functions $F$ and $G$ such that $H=G\om F$ is also an N-function
but
\[
  F^{-1}H^{-1}\not\approx G^{-1}.
\]
\end{theorem}

\begin{proof}
Let $A$ be the function from Lemma \ref{lem:A}, and define $F=A^{-1}$. Since
$A$ is increasing and concave, $F$ is increasing and convex. The limiting
relations in Lemma \ref{lem:A} give
\[
  \frac{F(t)}{t}\to0\quad(t\downarrow0),\qquad
  \frac{F(t)}{t}\to\infty\quad(t\to\infty),
\]
so $F$ is an N-function. Set
\[
  G(t)=t^2.
\]
Then $G$ is an N-function. Let
\[
  H(t)=(G\om F)(t)=\sup_{s\geq0}\{s^2t^2-F(s)\}.
\]

We first check that $H$ is an N-function. Since $A(u)/\sqrt u\to0$, writing
$s=A(u)$ gives
\[
  \frac{F(s)}{s^2}=\frac{u}{A(u)^2}\to\infty\qquad(s\to\infty).
\]
Hence for each fixed $t$, the quantity $s^2t^2-F(s)$ tends to $-\infty$ as
$s\to\infty$, so $H(t)<\infty$. As a supremum of convex increasing functions of
$t\geq0$, $H$ is a Young function.

Moreover, $H(t)/t\to\infty$ as $t\to\infty$, because fixing any $s_0>0$ gives
$H(t)\geq s_0^2t^2-F(s_0)$. For the behavior at zero, let $\varepsilon>0$.
Choose $R$ so large that $F(s)\geq s^2$ for $s\geq R$. If $0<t<1$, then for
$s\geq R$ we have $s^2t^2-F(s)\leq0$, while for $0\leq s\leq R$ we have
$s^2t^2-F(s)\leq R^2t^2$. Thus for sufficiently small $t$,
$H(t)\leq R^2t^2\leq\varepsilon t$. Therefore $H(t)/t\to0$ as $t\downarrow0$.
So $H$ is an N-function.

It remains to show failure of the inverse equivalence. Let $u_n,y_n$ be as in
Lemma \ref{lem:A}. Since $F(A(u_ny_n))=u_ny_n$, the definition of $H$ gives
for every $t>0$
\[
  H(t)\geq A(u_ny_n)^2t^2-u_ny_n.
\]
Consequently,
\[
  H^{-1}(u_n)\leq
  \frac{\sqrt{u_n+u_ny_n}}{A(u_ny_n)}.
\]
With the right-continuous inverse convention used for Orlicz functions, one may
first replace the numerator by $\sqrt{(1+\varepsilon)u_n+u_ny_n}$ and then let
$\varepsilon\downarrow0$; the displayed estimate is the resulting bound.
Using $G^{-1}(u_n)=\sqrt{u_n}$ and $F^{-1}=A$, we obtain
\[
  \frac{F^{-1}(u_n)H^{-1}(u_n)}{G^{-1}(u_n)}
  \leq
  \frac{A(u_n)}{A(u_ny_n)}\sqrt{1+y_n}
  \leq
  \frac{4\sqrt{1+y_n}}{y_n}\to0.
\]
Thus no constant $c>0$ can satisfy
$cG^{-1}\leq F^{-1}H^{-1}$, even along large arguments. Hence
$F^{-1}(G\om F)^{-1}\not\approx G^{-1}$.
\end{proof}

\section*{Verification and novelty notes}

The construction is deterministic and the proof is analytic. No computational
verification is needed. The only delicate point is the inverse construction in
Lemma \ref{lem:A}; it is included explicitly so that reviewers can check
convexity and the N-function limits directly.

Local indexes for the run were searched for arXiv:2410.01420, pointwise
multipliers, Calderon-Lozanovskii spaces, factorization, and related multiplier
packets. No exact prior packet was found. Web searches for the exact inverse
relation, for ``$G\om F$'' with N-functions, and for the source paper title
found the source arXiv page but no later answer to Problem 6.a.1.

\begin{thebibliography}{9}

\bibitem{KT}
T. Kiwerski and J. Tomaszewski,
\emph{A few last words on pointwise multipliers of Calderon-Lozanovskii spaces},
arXiv:2410.01420, 2024.

\end{thebibliography}

\end{document}
